\section{Falsifiability}

All results in this paper are formulated so as to admit explicit
falsification \emph{within} the Ontology of Continua (OC Core) and the
executable specification ETS.K\_EA.v1.1. Falsifiability is defined
strictly at the level of admissible observability constructs and logical
invariance properties; no external semantics, methods, or assumptions
are required.

\subsection{Falsifiability of decidability claims}

A decidability claim for a state predicate $\varphi$ asserts invariance
of $\varphi$ across all enterprise continua that are observationally
equivalent under the ETS-defined observation operator $\Pi$.

Such a claim is falsified if and only if there exist:
\begin{itemize}
\item two enterprise continua $K_{EA}^{(1)}$ and $K_{EA}^{(2)}$ that are
      observationally equivalent, i.e.\
      \[
      \Pi(\Omega_{EA}^{(1)}) = \Pi(\Omega_{EA}^{(2)});
      \]
\item an admissible state predicate $\varphi$ such that
      \[
      \varphi(\Omega_{EA}^{(1)}) \neq \varphi(\Omega_{EA}^{(2)}).
      \]
\end{itemize}

The existence of such a pair constitutes a direct counterexample and
invalidates the decidability claim. No probabilistic or algorithmic
considerations enter this criterion.

\subsection{Falsifiability of undecidability claims}

An undecidability claim asserts the \emph{existence} of at least one
observational equivalence class on which a predicate $\varphi$ is
non-invariant.

Such a claim is falsified if it can be shown that:
\begin{itemize}
\item for every pair of observationally equivalent enterprise continua,
      the predicate $\varphi$ agrees; equivalently,
\item there exists a mapping
      \[
      \tilde{\varphi} : \Sigma_{\mathrm{obs}} \to
      \{\text{true},\text{false}\}
      \]
      such that $\varphi = \tilde{\varphi} \circ \Pi$.
\end{itemize}

Either condition establishes invariance of $\varphi$ over all
observational equivalence classes and contradicts the undecidability
result.

\subsection{Falsifiability of STOP-only claims}

STOP-only decidability claims assert a one-sided invariance: STOP
membership is invariant under observational equivalence, while
non-STOP outcomes are not.

Such a claim is falsified by either of the following:
\begin{itemize}
\item the existence of observationally equivalent enterprise continua
      for which STOP holds in one case but not in the other;
\item a demonstration that at least one non-STOP phase outcome is also
      invariant across all observationally equivalent continua.
\end{itemize}

The first case violates STOP invariance directly. The second collapses
the STOP-only classification into full decidability.

\subsection{Observable-signature-level refutation}

All falsification conditions reduce to statements about observable
signatures in $\Sigma_{\mathrm{obs}}$ and their compatibility with
OC/ETS-admissible predicates. No falsification requires access to
unobservable internal structure beyond what is already constrained by
ETS observability.

In particular, a single admissible observable signature
$\sigma \in \Sigma_{\mathrm{obs}}$ that is compatible with mutually
inconsistent predicate values is sufficient to refute any claim of
decidability for that predicate.

\subsection{No meta-level escape}

No result in this paper can be falsified or defended by appeal to
improved measurement, richer data streams, probabilistic modeling, or
algorithmic inference \emph{within the same OC/ETS framework}. Any such
appeal would constitute a change of formal system and therefore lies
outside the scope of admissible falsification for the present results.

Within OC Core and ETS.K\_EA.v1.1, falsifiability is complete and
exhaustive at the level specified above.
